\documentclass[11pt,a4paper]{article}
\usepackage[margin=1in]{geometry}
\usepackage{amsmath,amssymb,amsfonts,amsthm}
\usepackage{booktabs}
\usepackage{hyperref}

\newtheorem{finding}{Finding}
\newtheorem{gap}{Gap}

\title{The $a_4$ Seeley--DeWitt Coefficient in DFD's $\alpha$ Derivation:\\
Cross-Reference Audit}

\author{Research Audit --- March 2026}
\date{}

\begin{document}
\maketitle

%=============================================================================
\section{Executive Summary}
%=============================================================================

After exhaustive search of all \texttt{.tex} files in both the v1.08 and v3.2
source trees, the standalone PDF ``Ab Initio Derivation of the Fine Structure
Constant from Density Field Dynamics'' (G.~Alcock, December 27, 2025), and
all supplementary notes, the following conclusion is reached:

\begin{center}
\fbox{\parbox{0.9\textwidth}{%
\textbf{No explicit formula for the $a_4$ Seeley--DeWitt coefficient on the
Toeplitz-truncated $\mathbb{C}P^2 \times S^3$ geometry appears anywhere in
the DFD corpus.}  The $\alpha$ derivation references $a_4$ as the source of
the gauge kinetic term but never writes out the curvature-invariant expression.
The derivation instead proceeds via an alternative route: Chern--Simons
partition function weights on $S^3$, combined with lattice Monte Carlo
verification and a Toeplitz-quantized spectral action with plateau cutoff.}}
\end{center}

%=============================================================================
\section{What the DFD Sources Actually Contain}
%=============================================================================

\subsection{Section \texttt{section\_alpha\_locked.tex}: The Spectral Action Framework}

This is the closest the corpus comes to an $a_4$-based derivation.  The key
elements present are:

\begin{enumerate}
\item \textbf{Operator choice (locked):} Connection Laplacian
$P = -g^{ij}\nabla_i\nabla_j$ on the internal bundle over
$X = \mathbb{C}P^2 \times S^3$ (Eq.~\ref{eq:conn-lap} of source, line 26).

\item \textbf{Bundle structure:} U(1) implemented via line bundle over
$\mathbb{C}P^2$ with curvature proportional to the K\"ahler form $\omega$,
trivial over $S^3$.  Stated property: $\nabla\omega = 0$, so derivative terms
in \emph{higher} Seeley--DeWitt coefficients vanish automatically (line 31).

\item \textbf{Reference to $a_4$:} ``The gauge-kinetic extraction from $a_4$
is unambiguous with this operator choice'' (line 34).  And: the plateau cutoff
``Preserves the leading $a_4$ gauge kinetic term'' (line 55).

\item \textbf{Plateau cutoff:} $S = \mathrm{Tr}\,f(P/\Lambda^2)$ with
$f^{(n)}(0) = 0$ for all $n \geq 1$, killing all negative moments
$f_{-2} = f_{-4} = \cdots = 0$.  This eliminates $a_6, a_8, \ldots$
contributions (lines 40--58).

\item \textbf{Toeplitz truncation:} At level $m = k+3$ on $\mathbb{C}P^1$,
giving $d = k+4 = 64$ for $k = 60$.  Spectral cutoff:
$\Lambda^3 = k \cdot (k+3)/(k+4)$ (line 79).

\item \textbf{The forced fork:} Regular-module microsector
$\mathcal{H}_F = M_d(\mathbb{C})$ vs.\ fermion-rep $\mathcal{H}_F = \mathbb{C}^d$.
Boost factor $\varepsilon_{\mathrm{adj}}^{(A)} = d^2/(d^2-1) = 4096/4095$
(line 114).

\item \textbf{Final numerical inputs:} $N_{\mathrm{species}} = 7$,
$\mathrm{Tr}(Y^2) = 10$, $g_F = 8$, $w = 7/80$ (Table, lines 197--201).

\item \textbf{Result:} $\alpha^{-1} = 137.03599985$ (residual $-0.006$ ppm).
\end{enumerate}

\textbf{What is missing:} The explicit expression for $a_4(P, X)$ in terms of
curvature invariants ($R$, $\mathrm{Ric}$, Riemann, bundle curvature $\Omega$)
evaluated on $\mathbb{C}P^2 \times S^3$.  The text asserts that $a_4$ gives
the gauge kinetic term but never writes the standard heat-kernel formula
\begin{equation}\label{eq:standard-a4}
a_4(P) = \frac{1}{(4\pi)^{d/2}} \int_X \mathrm{tr}\!\left(
\frac{1}{180}(R_{\mu\nu\rho\sigma}R^{\mu\nu\rho\sigma}
- R_{\mu\nu}R^{\mu\nu} + \tfrac{5}{2}R^2)
+ \frac{1}{12}\Omega_{\mu\nu}\Omega^{\mu\nu}
+ \cdots \right) \mathrm{dvol}
\end{equation}
nor evaluates it on the specific geometry.

\subsection{Appendix K (\texttt{appendix\_K.tex}): Microsector Physics}

Contains:
\begin{itemize}
\item Heat kernel on $S^3$: spectral expansion
$K(t; S^3) = \sum_{n=0}^{\infty}(n+1)^2 e^{-n(n+2)t/R^2}$ (line 41).
\item Spin$^c$ Dirac operator on $\mathbb{C}P^2$ identified with
$\sqrt{2}(\bar\partial + \bar\partial^*)$ (line 51).
\item Determination of $k_{\max} = 60$ via holomorphic Euler characteristic
$\chi(\mathbb{C}P^2, \mathcal{O}(9) \oplus \mathcal{O}^{\oplus 5}) = 55 + 5 = 60$
(line 68).
\item The $\alpha$ derivation itself via Chern--Simons weighted sum, \emph{not}
via explicit $a_4$ evaluation (lines 30--96).
\end{itemize}

\textbf{No $a_4$ formula appears in Appendix K.}

\subsection{Supplementary: \texttt{DFD\_Standard\_Model\_pasted\_note.tex}}

Contains:
\begin{itemize}
\item Topological coefficient $b = \dim(G) \times (\chi + 2\tau) = 12 \times 5 = 60$
(line 206).  This is related to heat-kernel counting of degrees of freedom but
is \emph{not} the $a_4$ coefficient.
\item Reference to Gilkey (1975) in the bibliography (line 441).
\item Bridge Lemma: $b = k_{\max} - h^{\vee} = 62 - 2 = 60$ connecting heat
kernel to Chern--Simons (line 67).
\end{itemize}

\subsection{The \texttt{section\_alpha\_relations.tex}: $C_{\mathrm{loop}} = 1/8$}

States: ``$C_{\mathrm{loop}} = 1/8$: Arises from the one-loop heat kernel
coefficient in the path integral.  This factor is \textbf{plausible from heat
kernel structure but requires explicit verification}'' (line 92, emphasis in
original).  Status listed as ``Plausible (B)'' (line 104).

This is the clearest admission that the heat-kernel (Seeley--DeWitt) computation
has not been carried through explicitly.

\subsection{The Ab Initio PDF}

The standalone paper presents the lattice Monte Carlo verification.  It states
(Section VI.C, page 21): ``the full derivation uses a Toeplitz-truncated
spectral action on $\mathbb{C}P^2 \times S^3$ and resolves a forced binary fork
\ldots\ Readers who wish to follow the complete derivation---including the
Toeplitz truncation and branch selection that make the derivation
`convention-locked'---should consult Appendix K of Ref.~[9].''

The PDF contains no $a_4$ formula.  The derivation route in the PDF is entirely
via Chern--Simons partition function weights.

%=============================================================================
\section{How the $\alpha$ Derivation Actually Works (Two Routes)}
%=============================================================================

DFD derives $\alpha$ via two complementary but distinct routes.  Neither
explicitly evaluates $a_4$.

\subsection{Route 1: Chern--Simons Weighted Sum (Primary)}

\begin{enumerate}
\item SU(2)$_k$ Chern--Simons partition function on $S^3$:
$Z_{SU(2)_k}(S^3) = \sqrt{2/(k+2)}\,\sin(\pi/(k+2))$.
\item Euclidean weight: $w(k) = |Z|^2 = \frac{2}{k+2}\sin^2\!\frac{\pi}{k+2}$.
\item Shifted level expectation:
$\langle k_{\mathrm{eff}} \rangle = \langle k+2 \rangle_{k_{\max}=60} = 3.7969 \approx 3.80$.
\item Dictionary: $\beta_{U(1)} = \langle k_{\mathrm{eff}} \rangle = 3.80$.
\item Lattice ratio from topology: $\beta_{SU(2)}/\beta_{U(1)} = (n_2/n_1) \times N_{\mathrm{gen}} = 2 \times 3 = 6$, giving $\beta_{SU(2)} = 22.80$.
\item Lattice MC at $(\beta_{U(1)}, \beta_{SU(2)}) = (3.80, 22.80)$ measures
stiffnesses $\kappa_{U(1)}, \kappa_{SU(2)}$.
\item Electroweak mixing:
$\alpha_W = \frac{(1/\kappa_{U(1)})(4/\kappa_{SU(2)})}{(1/\kappa_{U(1)}) + (4/\kappa_{SU(2)})} \cdot \frac{1}{4\pi}$.
\item Result: $\alpha_W \approx 0.007297 = 1/137.04$.
\end{enumerate}

\subsection{Route 2: Convention-Locked Spectral Action (Sub-ppm)}

\begin{enumerate}
\item Operator: connection Laplacian $P = -g^{ij}\nabla_i\nabla_j$ on internal bundle.
\item Spectral action: $S = \mathrm{Tr}\,f(P/\Lambda^2)$ with plateau cutoff
($f_{-2} = f_{-4} = \cdots = 0$).
\item \emph{Implicit} use of $a_4$: the gauge kinetic term is ``extracted from $a_4$''
but the coefficient is never written out.
\item Toeplitz quantization at $d = 64$; cutoff $\Lambda^3 = k(k+3)/(k+4)$.
\item Regular-module trace with boost $4096/4095$.
\item Hypercharge weighting: $w = N_{\mathrm{species}}/(g_F \cdot \mathrm{Tr}(Y^2)) = 7/80$.
\item Result: $\alpha^{-1} = 137.03599985$.
\end{enumerate}

The key observation is that Route 2 bypasses explicit evaluation of $a_4$ by:
(a) using the plateau cutoff to kill all higher $a_n$ terms, so only $a_4$
survives; (b) absorbing the geometric content of $a_4$ into the Toeplitz
truncation and trace normalization factors.

%=============================================================================
\section{The Standard $a_4$ Formula and What It Would Look Like on $\mathbb{C}P^2 \times S^3$}
%=============================================================================

For a Laplace-type operator $P = -(g^{ij}\nabla_i\nabla_j + E)$ on a
$d$-dimensional Riemannian manifold $X$, the standard Seeley--DeWitt--Gilkey
result (Gilkey 1975, Vassilevich 2003) gives:

\begin{equation}
a_4(P) = \frac{1}{(4\pi)^{d/2}} \frac{1}{360} \int_X \mathrm{tr}_V \!\Big(
60\,E\,R + 180\,E^2 + 30\,\Omega_{\mu\nu}\Omega^{\mu\nu}
+ (5R^2 - 2\,R_{\mu\nu}R^{\mu\nu} + 2\,R_{\mu\nu\rho\sigma}R^{\mu\nu\rho\sigma})
\cdot \mathrm{Id}_V
\Big)\,\mathrm{dvol}_g
\end{equation}

where $E$ is the endomorphism (zero for the pure connection Laplacian),
$\Omega_{\mu\nu}$ is the curvature of the bundle connection, and $\mathrm{tr}_V$
is the fiber trace.

\subsection{Evaluation on $\mathbb{C}P^2 \times S^3$ (What Should Be Done)}

For DFD's connection Laplacian ($E = 0$) with the U(1) line bundle twisted by
$\omega$ on $\mathbb{C}P^2$ (trivial on $S^3$):

\begin{enumerate}
\item $\mathbb{C}P^2$ is K\"ahler--Einstein with $R_{\mu\nu} = 6g_{\mu\nu}$
(for unit Fubini--Study metric), $R = 24$,
$R_{\mu\nu\rho\sigma}R^{\mu\nu\rho\sigma} = 24$.

\item $S^3$ has constant positive curvature: $R_{S^3} = 6/r^2$,
$R_{\mu\nu}R^{\mu\nu} = 12/r^4$,
$R_{\mu\nu\rho\sigma}R^{\mu\nu\rho\sigma} = 12/r^4$.

\item The bundle curvature $\Omega = i\omega$ on $\mathbb{C}P^2$ (proportional
to K\"ahler form), zero on $S^3$.  Hence
$\mathrm{tr}(\Omega_{\mu\nu}\Omega^{\mu\nu})$ reduces to a topological
invariant on $\mathbb{C}P^2$.

\item Product geometry: $a_4(X) = a_0(\mathbb{C}P^2) \cdot a_4(S^3)
+ a_2(\mathbb{C}P^2) \cdot a_2(S^3) + a_4(\mathbb{C}P^2) \cdot a_0(S^3)$
by the product formula for heat-kernel coefficients.

\item The gauge kinetic term emerges from the $\Omega_{\mu\nu}\Omega^{\mu\nu}$
piece, which on $\mathbb{C}P^2$ gives $\propto \int_{\mathbb{C}P^2} \omega \wedge \omega
= \mathrm{Vol}(\mathbb{C}P^2)/2 \cdot c_1^2$.
\end{enumerate}

The \textbf{Toeplitz truncation} at level $d = 64$ then replaces the
continuous $\mathbb{C}P^2$ integral with a finite-dimensional matrix trace,
modifying the overall normalization by the factor $(k+3)/(k+4) = 63/64$ and
introducing the boost factor $d^2/(d^2-1) = 4096/4095$.

\textbf{This explicit evaluation has not been published in any DFD source.}

%=============================================================================
\section{Identified Gaps}
%=============================================================================

\begin{gap}[No explicit $a_4$ formula]
The DFD corpus refers to the gauge kinetic term being ``extracted from $a_4$''
but never writes the Seeley--DeWitt--Gilkey formula evaluated on
$\mathbb{C}P^2 \times S^3$ with the specific bundle data.
\end{gap}

\begin{gap}[$C_{\mathrm{loop}} = 1/8$ unverified]
The self-coupling relation $k_a = 3/(8\alpha)$ depends on a factor
$C_{\mathrm{loop}} = 1/8$ that is described as ``plausible from heat kernel
structure but requires explicit verification'' (section\_alpha\_relations.tex,
line 92).  This factor should emerge from the $a_4$ computation.
\end{gap}

\begin{gap}[Product-formula heat-kernel decomposition absent]
The product formula $a_4(X) = \sum_{p+q=4} a_p(\mathbb{C}P^2) \cdot a_q(S^3)$
is never invoked or evaluated.  This would make the geometric content of the
$\alpha$ derivation fully transparent.
\end{gap}

\begin{gap}[Toeplitz modification of $a_4$ not derived]
How exactly the Toeplitz truncation at level $d = 64$ modifies the continuous
$a_4$ integral into a finite matrix trace is stated as a result (the
$(k+3)/(k+4)$ factor and boost) but not derived from first principles.
\end{gap}

\begin{gap}[Connection to Chamseddine--Connes spectral action not explicit]
The DFD spectral action framework closely parallels noncommutative geometry
(Chamseddine--Connes 1996), where $a_4$ determines the gauge coupling via
$g^{-2} \propto f_0 \cdot a_4$.  The DFD sources never make this connection
explicit or compare their setup to the NCG spectral action.
\end{gap}

%=============================================================================
\section{What \emph{Is} Present (Positive Findings)}
%=============================================================================

\begin{finding}[Seeley--DeWitt referenced by name]
The term ``Seeley--DeWitt'' appears in \texttt{section\_alpha\_locked.tex}
(line 31): ``derivative terms in higher Seeley--DeWitt coefficients vanish
automatically'' due to $\nabla\omega = 0$.
\end{finding}

\begin{finding}[$a_4$ referenced as source of gauge kinetic term]
Lines 34 and 55 of \texttt{section\_alpha\_locked.tex} explicitly state that
the gauge kinetic term is ``extracted from $a_4$'' and that the plateau cutoff
``preserves the leading $a_4$ gauge kinetic term.''
\end{finding}

\begin{finding}[Heat kernel on $S^3$ given explicitly]
Appendix K provides the spectral expansion of the heat kernel on $S^3$
with eigenvalues $\lambda_n = n(n+2)/R^2$ and degeneracies $(n+1)^2$.
\end{finding}

\begin{finding}[Gilkey cited]
The bibliography in \texttt{DFD\_Standard\_Model\_pasted\_note.tex} cites
Gilkey (1975) --- the foundational reference for Seeley--DeWitt coefficients.
\end{finding}

\begin{finding}[Topological coefficient $b = 60$ related to heat kernel]
The formula $b = \dim(G) \times (\chi + 2\tau) = 12 \times 5 = 60$ in the
supplementary note is a heat-kernel--type invariant (related to the $a_4$
coefficient of the Hodge Laplacian on 1-forms twisted by the adjoint bundle).
\end{finding}

\begin{finding}[Full numerical chain for $\alpha^{-1}$ is present]
Despite the absence of an explicit $a_4$ formula, the complete numerical
derivation chain yielding $\alpha^{-1} = 137.03599985$ is fully documented
in \texttt{section\_alpha\_locked.tex} with every input identified and locked.
\end{finding}

%=============================================================================
\section{Recommendation}
%=============================================================================

The DFD corpus would benefit from a dedicated subsection (either in
\texttt{section\_alpha\_locked.tex} or a new appendix) that:

\begin{enumerate}
\item Writes the standard $a_4$ Seeley--DeWitt--Gilkey formula for the connection
Laplacian.

\item Evaluates it on $\mathbb{C}P^2 \times S^3$ using the product formula,
inserting the known curvature invariants of both factors and the U(1) bundle
curvature $\Omega = i\omega$.

\item Shows how the Toeplitz truncation at level $d = 64$ replaces the
continuous integral with a finite matrix trace, deriving the $(k+3)/(k+4)$
factor and the boost $d^2/(d^2-1)$ from the Toeplitz symbol calculus.

\item Demonstrates that the plateau cutoff isolates exactly the $a_4$ term,
so the result is $g^{-2} \propto f_0 \cdot a_4^{\mathrm{Toeplitz}}$, and
derives the final $\alpha^{-1}$ from this expression.

\item Verifies that $C_{\mathrm{loop}} = 1/8$ emerges from this computation,
upgrading its status from ``Plausible (B)'' to ``Derived (A).''
\end{enumerate}

This would close the gap between the \emph{assertion} that $a_4$ determines
the gauge coupling and the \emph{demonstration} thereof, making the
convention-locked derivation fully self-contained.

%=============================================================================
\section{Source File Locations}
%=============================================================================

\begin{center}
\small
\begin{tabular}{p{0.55\textwidth}p{0.35\textwidth}}
\toprule
\textbf{File} & \textbf{Relevant Content} \\
\midrule
\texttt{section\_alpha\_locked.tex} & Spectral action, $a_4$ references, Toeplitz truncation, forced fork, sub-ppm result \\
\texttt{appendix\_K.tex} & Heat kernel on $S^3$, Bridge Lemma, CS-based $\alpha$ derivation \\
\texttt{section\_alpha\_relations.tex} & $C_{\mathrm{loop}} = 1/8$ (status: Plausible B) \\
\texttt{DFD\_Standard\_Model\_pasted\_note.tex} & $b = 60$ from heat kernel, Gilkey citation \\
\texttt{appendix\_Z\_complete\_params.tex} & Weinberg angle from partition, spectral/proper-time regulator \\
\texttt{PRL\_SM\_from\_Topology.tex} & Summary of parameter derivation \\
\texttt{section\_gauge\_coupling.tex} & Gauge coupling variation, $\beta$-function \\
\texttt{DFD\_Unified\_Review.tex} & Abstract mentions ``microsector spectral action on $\mathbb{C}P^2 \times S^3$'' \\
\texttt{Ab\_Initio\ldots.pdf} & Lattice MC verification; no $a_4$ formula \\
\bottomrule
\end{tabular}
\end{center}

\end{document}
